
\newpage
\begin{figure}[htbp]
\centering
\makebox[\textwidth][c]{%
\begin{tikzpicture}[
    >=Stealth,
    node distance=1.5cm,
    box/.style={rectangle, draw=border, rounded corners, minimum width=2.5cm, minimum height=0.8cm, align=center, fill=paleGray},
    llmbox/.style={rectangle, draw=darkCyan, thick, rounded corners=8pt, minimum width=3cm, minimum height=1.2cm, align=center, fill=paleCyan, font=\bfseries},
    probbox/.style={rectangle, draw=border, rounded corners, minimum width=4cm, minimum height=1.5cm, align=center, fill=paleGray},
    tokenbox/.style={rectangle, draw=border, minimum width=0.6cm, minimum height=0.6cm, fill=paleCyan, font=\small\ttfamily},
    newtoken/.style={rectangle, draw=darkGray, minimum width=0.6cm, minimum height=0.6cm, fill=coolGray, font=\small\ttfamily},
    arrow/.style={->, thick, >=Stealth, color=darkGray},
    looparrow/.style={->, very thick, >=Stealth, color=darkCyan}
]

% === 左侧：输入序列 ===
\node[anchor=west] (inputlabel) at (-6.5, 0) {\textbf{\ifdefined\ENGLISH Current sequence $x$\else 当前序列 $x$\fi}};

% Token序列
\ifdefined\ENGLISH
  \node[tokenbox] (t1) at (-6, -1) {The};
  \node[tokenbox, right=0.1cm of t1] (t2) {cat};
  \node[tokenbox, right=0.1cm of t2] (t3) {sat};
  \node[tokenbox, right=0.1cm of t3] (t4) {on};
  \node[newtoken, right=0.1cm of t4] (t5) {?};
\else
  \node[tokenbox] (t1) at (-6, -1) {今};
  \node[tokenbox, right=0.1cm of t1] (t2) {天};
  \node[tokenbox, right=0.1cm of t2] (t3) {天};
  \node[tokenbox, right=0.1cm of t3] (t4) {气};
  \node[newtoken, right=0.1cm of t4] (t5) {?};
\fi

% === 中间：LLM ===
\node[llmbox, minimum width=2.2cm, minimum height=1cm] (llm) at (-0.5, -1) {LLM\\$P_\theta(y|x)$};

% === 右侧：概率分布 ===
\node[anchor=west] (problabel) at (0.8, 0) {\textbf{\ifdefined\ENGLISH Next-token probability distribution\else 下一个Token的概率分布\fi}};

% 概率条形图
\begin{scope}[shift={(2.5, -3.3)}]
    \draw[rounded corners, fill=paleGray, draw=border] (-1.2, -1.5) rectangle (2.1, 1.5);    
    \draw[fill=accentCyan] (-0.5, 1) rectangle (1.5, 1.3);
    \node[left, font=\small\ttfamily] at (-0.5, 1.15) {\ifdefined\ENGLISH the\else 很\fi};
    \node[right, font=\scriptsize] at (1.5, 1.15) {0.35};

    \draw[fill=cyan] (-0.5, 0.5) rectangle (1.2, 0.8);
    \node[left, font=\small\ttfamily] at (-0.5, 0.65) {\ifdefined\ENGLISH a\else 真\fi};
    \node[right, font=\scriptsize] at (1.2, 0.65) {0.28};

    \draw[fill=lightCyan] (-0.5, 0) rectangle (0.8, 0.3);
    \node[left, font=\small\ttfamily] at (-0.5, 0.15) {\ifdefined\ENGLISH my\else 不\fi};
    \node[right, font=\scriptsize] at (0.8, 0.15) {0.18};

    \draw[fill=midCyan] (-0.5, -0.5) rectangle (0.5, -0.2);
    \node[left, font=\small\ttfamily] at (-0.5, -0.35) {\ifdefined\ENGLISH it\else 挺\fi};
    \node[right, font=\scriptsize] at (0.5, -0.35) {0.12};
    
    \draw[fill=softCyan] (-0.5, -1) rectangle (0.3, -0.7);
    \node[left, font=\small\ttfamily] at (-0.5, -0.85) {...};
    \node[right, font=\scriptsize] at (0.3, -0.85) {0.07};
\end{scope}

% === 底部：采样结果 ===
\node[align=center] (samplelabel) at (-1.5, -5.5) {\textbf{\ifdefined\ENGLISH Sample\else 采样\fi} $y \sim P_\theta(y|x)$};
\node[newtoken, below=0.2cm of samplelabel] (sampled) {\ifdefined\ENGLISH the\else 很\fi};

% === 箭头连接 ===
\draw[arrow] (t5.east) -- (llm.west);
\draw[arrow] (llm.east) -- (3.2, -1)  -- (3.2, -1.8);
\draw[arrow] (3.2, -4.8) -- (3.2, -6.3) -- (sampled.east);

% === 循环箭头 ===
\draw[looparrow, rounded corners=15pt] 
    (sampled.west) -- (-7, -6.3) -- (-7, -1) -- (t1.west);

\node[fill=paleCyan, draw=border, rounded corners, font=\small, align=center] at (-7, -3.8) {\ifdefined\ENGLISH Append to\\sequence\else 追加到\\序列末尾\fi};

% === 下方：时间线视图 ===
\begin{scope}[shift={(-2.5, -8)}]
    \node[anchor=west, font=\bfseries] at (-5, 1) {\ifdefined\ENGLISH Unrolled time view:\else 时间展开视图：\fi};

    \node[font=\small] at (-4.5, 0) {$t=0$};
    \ifdefined\ENGLISH
      \node[tokenbox, scale=0.8] at (-4.5, -0.7) {The};
      \node[tokenbox, scale=0.8] at (-3.9, -0.7) {cat};
      \node[tokenbox, scale=0.8] at (-3.3, -0.7) {sat};
      \node[tokenbox, scale=0.8] at (-2.7, -0.7) {on};
    \else
      \node[tokenbox, scale=0.8] at (-4.5, -0.7) {今};
      \node[tokenbox, scale=0.8] at (-3.9, -0.7) {天};
      \node[tokenbox, scale=0.8] at (-3.3, -0.7) {天};
      \node[tokenbox, scale=0.8] at (-2.7, -0.7) {气};
    \fi

    \draw[arrow, color=border] (-2.1, -0.7) -- (-1.5, -0.7);

    \node[font=\small] at (-0.5, 0) {$t=1$};
    \ifdefined\ENGLISH
      \node[tokenbox, scale=0.8] at (-1.1, -0.7) {The};
      \node[tokenbox, scale=0.8] at (-0.5, -0.7) {cat};
      \node[tokenbox, scale=0.8] at (0.1, -0.7) {sat};
      \node[tokenbox, scale=0.8] at (0.7, -0.7) {on};
      \node[newtoken, scale=0.8] at (1.3, -0.7) {the};
    \else
      \node[tokenbox, scale=0.8] at (-1.1, -0.7) {今};
      \node[tokenbox, scale=0.8] at (-0.5, -0.7) {天};
      \node[tokenbox, scale=0.8] at (0.1, -0.7) {天};
      \node[tokenbox, scale=0.8] at (0.7, -0.7) {气};
      \node[newtoken, scale=0.8] at (1.3, -0.7) {很};
    \fi

    \draw[arrow, color=border] (1.9, -0.7) -- (2.5, -0.7);

    \node[font=\small] at (3.7, 0) {$t=2$};
    \ifdefined\ENGLISH
      \node[tokenbox, scale=0.8] at (2.7, -0.7) {The};
      \node[tokenbox, scale=0.8] at (3.3, -0.7) {cat};
      \node[tokenbox, scale=0.8] at (3.9, -0.7) {sat};
      \node[tokenbox, scale=0.8] at (4.5, -0.7) {on};
      \node[tokenbox, scale=0.8] at (5.1, -0.7) {the};
      \node[newtoken, scale=0.8] at (5.7, -0.7) {mat};
    \else
      \node[tokenbox, scale=0.8] at (2.7, -0.7) {今};
      \node[tokenbox, scale=0.8] at (3.3, -0.7) {天};
      \node[tokenbox, scale=0.8] at (3.9, -0.7) {天};
      \node[tokenbox, scale=0.8] at (4.5, -0.7) {气};
      \node[tokenbox, scale=0.8] at (5.1, -0.7) {很};
      \node[newtoken, scale=0.8] at (5.7, -0.7) {好};
    \fi
    
    \node at (6.5, -0.7) {$\cdots$};
    
    \node[font=\small, color=darkGray, align=center] at (0.5, -1.8) {\ifdefined\ENGLISH Each step: input sequence $\rightarrow$ predict distribution $\rightarrow$ sample token $\rightarrow$ append\else 每一步：输入当前序列 $\rightarrow$ 预测概率分布 $\rightarrow$ 采样新Token $\rightarrow$ 追加到序列\fi};
\end{scope}

\end{tikzpicture}
}
\caption{\ifdefined\ENGLISH Given current sequence $x$, the model outputs a probability distribution $P_\theta(y|x)$ for the next token; after sampling, the new token is appended and the process repeats until the task is complete\else 给定当前序列$x$，模型输出下一个Token的概率分布$P_\theta(y|x)$，采样得到新Token后将其追加到序列末尾，循环往复直至完成任务\fi}
\label{fig:autoregressive}
\end{figure}
